\section{Expertise, Kreativität, Intelligenz}
\slide{Zusammenhang Expertise \& Intelligenz}{
    \begin{itemize}
        \item \b{Abgrenzung:} Expertise ist nicht gleich Intelligenz und umgekehrt.
        \item \b{Hypothese:} Expertise korreliert eher mit kristalliner Intelligenz (Wissen), Kreativität eher mit fluider Intelligenz (Flexibilität).
    \end{itemize}
}
\slide{Expertise}{
    \begin{itemize}
        \item Definition: Außergewöhnliche Fähigkeiten und Kenntnisse in einem spezifischen Bereich (domänen-spezifisch).
        \item Erwerb: Expertise entsteht primär durch intensive Übung und Bildung (Bereichswissen), nicht allein durch Talent.
        \item Wissensorganisation: Experten verfügen über viel deklaratives (Fakten) und prozedurales Wissen (Strategien) sowie ausgeprägte Metakognition.
        \begin{itemize}
            \item Tiefenverarbeitung: Experten sortieren Probleme nach zugrundeliegenden Lösungsprinzipien (Tiefenstruktur), Novizen nach oberflächlichen Merkmalen (z. B. "da ist eine schiefe Ebene").
        \end{itemize}
        \item Gedächtnisleistung (Schachbrett-Experiment): Experten erinnern sinnvolle Muster deutlich besser als Novizen (Chunking durch Fachwissen). Bei zufälligen Anordnungen verschwindet der Vorteil, da das Vorwissen nicht greift.
        \item Strategien: Experten nutzen effiziente Strategien wie Forward-chained-reasoning (sicheres Schließen vom Start zum Ziel) statt Backward-chained-reasoning (hypothetisches Rückwärtsschließen) und haben viele Prozesse automatisiert.
    \end{itemize}
}
\slide{Kreativität}{
    \begin{itemize}
        \item Definition: Das Erzeugen von etwas, das sowohl neu (originell) als auch nützlich (adäquat) ist.
        \item \b{Divergentes Denken:} Offene, unsystematische Suche nach vielen verschiedenen Lösungen (Brainstorming), im Gegensatz zu konvergentem Denken (eine richtige Lösung finden).
        \item \b{Investitions-Theorie (4 P's):} Kreativität wird beeinflusst durch Person (Persönlichkeit), Process (Denkweise), Product (Ergebnis) und Place (förderndes Umfeld).
        \item Inkubation: Eine Ruhephase, in der ein Problem unterbewusst weiterverarbeitet wird ("latente Wirkung"), was oft zu plötzlichen Einsichten führt.
    \end{itemize}
}
\slide{Intelligenz}{
    \begin{itemize}
        \item Definition: Allgemeine kognitive Leistungsfähigkeit; in der Wissenschaft als relativ konstante Persönlichkeitseigenschaft verstanden.
        \item \b{Faktorenmodell (Cattell):}
        \begin{itemize}
            \item item Fluide Intelligenz (G-Faktor): Angeborene, biologische Basisprozesse (z. B. Verarbeitungsgeschwindigkeit, Arbeitsgedächtnis); nimmt im Alter ab.
            \item Kristalline Intelligenz (S-Faktor): Erworbenes, kulturabhängiges Wissen und Wortschatz; bleibt über das Leben stabil oder nimmt zu.
        \end{itemize}
        \item Messung (IQ): Vergleich eines individuellen Rohwertes mit einer Normstichprobe (Mittelwert = 100, Standardabweichung = 15)
        \item Cannabis-Studie: Untersuchung des Zusammenhangs zwischen Cannabiskonsum und IQ-Entwicklung.
        \begin{itemize}
            \item item Ergebnis: IQ sank besonders stark, wenn der Konsum bereits in der Jugend begann und häufig stattfand.
            \item Einschränkung: Da es keine kontrollierte experimentelle Studie war, sind keine zwingenden Kausalschlüsse möglich (Korrelation vs. Kausalität).
        \end{itemize}
    \end{itemize}
}
